\documentclass[11pt]{article}
\usepackage[utf8]{inputenc}
\usepackage[spanish]{babel}
\usepackage[margin=2.5cm]{geometry}


\title{\textbf{Proyecto 1: Reconocimiento de Figuras} \\ \large \textbf{Reporte Técnico}}
\author{
    Integrantes:
}
\date{Modelado y Programación \\ \today}

\begin{document}

\maketitle

\section{Definición del problema}

Dada una imagen en formato $.bmp$ que contiene una o más figuras geométricas
de colores sólidos sobre un fondo también de color sólido, el programa debe reportar para cada
figura presente en la imagen:

\begin{itemize}
    \item Su color, en formato hexadecimal.
    \item Su clasificación en:
    \begin{itemize}
        \item \textbf{C} Cuadriláteros.
        \item \textbf{T} Triángulos.
        \item \textbf{O} Círculos.
        \item \textbf{X} Otros (cualquier otra figura que no sea de las categorías anteriores).
    \end{itemize}
\end{itemize}

\textbf{Pequeño ejemplo de salida:}\\ Forma: C \\ Color: Rojo %Solo es idea, falta corregir y mostrar buen ejemplo de la salida que da el programa.

\section{Arsenal}

\subsection{Lenguaje de programación}

\begin{itemize}
    \item \textbf{Python:} Nos decidimos por $Python$ como lenguaje de desarrollo porque cuenta con un buen ecosistema de bibliotecas para procesamiento de imágenes.
    Además, tiene versatilidad y una sintaxis simple, lo cual facilita el trabajo colaborativo
    y la familiarización con el lenguaje, otro punto a favor es que evita tener que preocuparse por cosas como punteros o memoria, presentes en otros lenguajes. 
\end{itemize}

\subsection{Bibliotecas}

\begin{itemize}
    \item \textbf{OpenCV:} Se eligió OpenCV por cubrir tanto la lectura de imagénes 
    como las operaciones de detección de contornos, formas, segmentación, visión por computadora, etc.
    Usar esta biblioteca agiliza el trabajo al evitar reimplementar estos algoritmos desde cero,
    a la vez que reduce un poco la probabilidad de caer en errores de las operaciones.
\end{itemize}

Vale la pena mencionar que, se consideraron tres bibliotecas para el manejo de imágenes: $Pillow$, útil para operaciones de lectura, manipulación de píxeles, 
colores y dimensiones pero limitada en análisis geométrico, $scikit-image$, orientada a procesamiento científico 
con segmentación, regiones además de propiedades geométricas y $OpenCV$, especializada en visión por computadora con soporte directo para detección de contornos 
y análisis de formas. Al tener una mejor oferta que se adaptaba a las necesidades del proyecto se optó por esta última opción.

\section{Análisis del problema}


\subsection{Supuestos}

Mismos supuestos dados en el detallamiento del proyecto. 
\begin{itemize}
    \item El fondo de la imagen es de un solo color, uniforme.
    \item Cada figura tiene un color sólido, distinto entre figuras y distinto del fondo.
    \item No hay anti-aliasing ni colores intermedios en los bordes de las figuras.
    \item Las figuras no se traslapan entre sí.
    \item Tamaño rotación y posición de las figuras son arbitrarios.
\end{itemize}

\subsection{Casos límite identificados}

\begin{itemize}
    \item La imagen no cumple con las especificaciones del problema (sea una fotografía/imagen compleja con elementos distintos a las figuras esperadas). 
    \textrightarrow{ Debe ser detectado por el programa y responder con un mensaje de error claro ('Imagen inválida') en vez de falle o dar una clasificación sin sentido.} 
    \item La imagen no esta en el formato adecuado (no es $.bmp$).
    \textrightarrow{ De igual manera se debe detectar ya que no se puede admitir esa entrada y responder con un mensaje ('Formato inválido').}
    \item Imagen sin ninguna figura, solo un fondo de color sólido. 
\end{itemize}

\subsection{Requisitos funcionales}

\begin{itemize}
    \item Lectura de imagénes mediante la dirección de un archivo $.bmp$
    \item Detección de figuras (C, T, O, X)
    \item Detección de colores (formato hexadecimal).
    \item Validación de entrada (formato de archivo e imagen)
    \item Algoritmo de búsqueda/clasificación de figuras
    \item Reporte de resultados al usuario: De forma clara y breve.
\end{itemize}

\subsection{Requisitos NO funcionales}

\begin{itemize}
    \item Rendimiento: El programa debe procesar una imagen de tamaño razonable en un tiempo aceptable.
    \item Usabilidad: El programa regresa resultados comprensibles para el usuario.
    \item Mantenibilidad:El código debe estar organizado en secciones con tareas bien definidas.
    \item Robustez: El programa debe manejar entradas inválidas sin terminar de manera inesperada.
    \item Portabilidad: El programa debe poder ejecutarse en los entornos definidos del proyecto.
    \item Documentación del código: El código debe tener documentación suficiente para facilitar la comprensión y futuras modificaciones.
\end{itemize}

\section{Selección de la mejor alternativa} 
%Algoritmo usado para la deteccion y clasificaciomn de figuras (¿porque lo preferimos?)

\section{Diagrama de flujo o pseudocódigo}

\end{document}